\documentclass[12pt]{article}
\usepackage[utf8]{inputenc}
\usepackage[T2A]{fontenc}
\usepackage[russian]{babel}
\usepackage{url}

\usepackage{amsmath, amssymb}
\usepackage{physics}
\usepackage{geometry}
\usepackage{graphicx}
\usepackage{setspace}
\geometry{margin=2.5cm}

\title{Physics-Informed Neural Networks:\\
интуитивное и теоретическое введение}
\author{}
\date{}

\begin{document}

\section{Мотивация}

Во многих задачах механики сплошных сред и геофизики поведение системы описывается уравнениями в частных производных. Примерами таких задач являются распространение упругих и термоупругих волн, теплоперенос, диффузия и динамика деформируемых сред.

Классические численные методы, такие как метод конечных элементов или конечных разностей, требуют построения пространственной сетки, что усложняет моделирование сложных геометрий и неоднородных сред. Кроме того, при наличии ограниченного объёма экспериментальных данных такие методы не всегда позволяют эффективно учитывать измерения.

Physics-Informed Neural Networks (PINN) представляют собой альтернативный подход, в котором дифференциальные уравнения задачи напрямую включаются в процесс обучения нейросети.


\section{Интуитивное объяснение PINN}

Обычная нейросеть учится по принципу:
\begin{quote}
``Вот входные данные, вот правильный ответ — постарайся их сопоставить''.
\end{quote}

PINN добавляет к этому принципу ещё одно требование:
\begin{quote}
``Твой ответ должен не только совпадать с данными, но и подчиняться законам физики''.
\end{quote}

Идея заключается в том, что нейросеть аппроксимирует неизвестные физические поля (например, смещения и температуру), а затем проверяется, насколько хорошо эти аппроксимации удовлетворяют дифференциальным уравнениям задачи.

Если физические уравнения нарушаются, нейросеть получает штраф в функции потерь и корректирует свои параметры.


\section{Математическая постановка задачи в контексте PINN}

Пусть задача описывается системой уравнений в частных производных:
\begin{equation}
\mathcal{F}\bigl(u(x,t), T(x,t)\bigr) = 0,
\end{equation}
где $u(x,t)$ — вектор перемещений, а $T(x,t)$ — температурное поле.

В методе PINN вводится нейросетевая аппроксимация
\begin{equation}
(u(x,t), T(x,t)) \approx \mathcal{N}_\theta(x,t),
\end{equation}
где $\mathcal{N}_\theta$ — нейросеть с параметрами $\theta$.

Важно, что нейросеть аппроксимирует не значения на сетке, а непрерывную функцию во всей области определения.

\section{Роль автоматического дифференцирования}

Одним из ключевых элементов PINN является автоматическое дифференцирование. Оно позволяет точно вычислять производные выхода нейросети по входным переменным:
\begin{equation}
\frac{\partial u}{\partial t}, \quad
\frac{\partial^2 u}{\partial x^2}, \quad
\nabla^2 T, \quad \text{и т.д.}
\end{equation}

Таким образом, дифференциальные операторы применяются не к численным данным, а непосредственно к нейросетевой аппроксимации, что исключает ошибки дискретизации.

\section{Функция потерь в PINN}

Функция потерь PINN состоит из нескольких частей:
\begin{equation}
\mathcal{L} =
\mathcal{L}_{\text{phys}} +
\mathcal{L}_{\text{BC}} +
\mathcal{L}_{\text{IC}} +
\mathcal{L}_{\text{data}},
\end{equation}
где:
\begin{itemize}
    \item $\mathcal{L}_{\text{phys}}$ — штраф за нарушение дифференциальных уравнений;
    \item $\mathcal{L}_{\text{BC}}$ — выполнение граничных условий;
    \item $\mathcal{L}_{\text{IC}}$ — выполнение начальных условий;
    \item $\mathcal{L}_{\text{data}}$ — соответствие экспериментальным данным (если они есть).
\end{itemize}

Физическая часть функции потерь имеет вид
\begin{equation}
\mathcal{L}_{\text{phys}} =
\left\| \mathcal{F}\bigl(\mathcal{N}_\theta(x,t)\bigr) \right\|^2,
\end{equation}
то есть нейросеть наказывается за несоблюдение уравнений движения и теплопереноса.

\section{Почему PINN особенно хорошо подходят для задач термоупругости}

Задачи термоупругости обладают рядом особенностей:
\begin{itemize}
    \item система уравнений является связанной;
    \item присутствуют пространственные и временные производные высоких порядков;
    \item параметры среды могут быть неоднородными;
    \item экспериментальные данные, как правило, ограничены.
\end{itemize}

PINN естественным образом учитывают все эти факторы, так как физические уравнения являются основным источником информации для обучения.

Кроме того, в PINN отсутствует необходимость в построении сетки, что упрощает моделирование сложных геологических сред.

\section{Концептуальный алгоритм решения задачи с помощью PINN}

Решение задачи термоупругих волн с использованием PINN включает следующие шаги:
\begin{enumerate}
    \item Задание нейросетевой аппроксимации искомых полей.
    \item Генерация коллокационных точек в пространственно-временной области.
    \item Вычисление производных через автоматическое дифференцирование.
    \item Формирование функции потерь на основе физических уравнений.
    \item Оптимизация параметров нейросети.
\end{enumerate}

В результате получается непрерывное приближение решения задачи во всей области.

\section{Заключение}

Physics-Informed Neural Networks представляют собой гибридный подход, объединяющий методы машинного обучения и классическую математическую физику. В задачах термоупругой волновой динамики PINN позволяют эффективно учитывать физические законы, работать с ограниченным объёмом данных и моделировать сложные неоднородные среды.

\end{document}
